\documentclass[11pt]{article}

\usepackage[letterpaper,margin=0.78in]{geometry}
\usepackage[T1]{fontenc}
\usepackage[utf8]{inputenc}
\usepackage{lmodern}
\usepackage{microtype}
\usepackage[table]{xcolor}
\usepackage{booktabs}
\usepackage{tabularx}
\usepackage{longtable}
\usepackage{array}
\usepackage{enumitem}
\usepackage{hyperref}
\usepackage{xurl}
\usepackage{fancyhdr}
\usepackage{titlesec}
\usepackage[most]{tcolorbox}

\definecolor{accent}{HTML}{1F4E79}
\definecolor{lightaccent}{HTML}{EAF2F8}
\definecolor{midgray}{HTML}{666666}
\definecolor{lightgray}{HTML}{F5F7FA}
\definecolor{warn}{HTML}{8A4B08}
\definecolor{success}{HTML}{2E7D32}

\hypersetup{
  colorlinks=true,
  linkcolor=accent,
  urlcolor=accent,
  citecolor=accent,
  pdftitle={JAN-PRO Franchise Sales Director Meeting Prep},
  pdfauthor={Prepared with ChatGPT},
  pdfsubject={JAN-PRO franchise candidate meeting preparation},
  pdfkeywords={JAN-PRO, franchise, commercial cleaning, Franchise Disclosure Document, FDD, SBA, FTC, due diligence}
}

\setlength{\headheight}{14pt}
\pagestyle{fancy}
\fancyhf{}
\lhead{JAN-PRO Franchise Meeting Prep}
\rhead{Candidate Positioning Handout}
\cfoot{\thepage}
\renewcommand{\headrulewidth}{0.4pt}

\titleformat{\section}{\Large\bfseries\color{accent}}{}{0em}{}[\titlerule]
\titleformat{\subsection}{\large\bfseries\color{accent}}{}{0em}{}
\titleformat{\subsubsection}{\normalsize\bfseries\color{accent}}{}{0em}{}
\setlist[itemize]{leftmargin=1.35em,itemsep=0.2em,topsep=0.25em}
\setlist[enumerate]{leftmargin=1.45em,itemsep=0.2em,topsep=0.25em}
\renewcommand{\arraystretch}{1.22}
\emergencystretch=2em
\sloppy

\newtcolorbox{keybox}[1]{
  colback=lightaccent,
  colframe=accent,
  title=\textbf{#1},
  arc=2mm,
  boxrule=0.6pt,
  left=1.1mm,
  right=1.1mm,
  top=0.9mm,
  bottom=0.9mm
}

\newtcolorbox{scriptbox}[1]{
  colback=lightgray,
  colframe=midgray,
  title=\textbf{#1},
  arc=1.5mm,
  boxrule=0.45pt,
  left=1.1mm,
  right=1.1mm,
  top=0.9mm,
  bottom=0.9mm
}

\newtcolorbox{warningbox}[1]{
  colback=orange!7,
  colframe=warn,
  title=\textbf{#1},
  arc=1.5mm,
  boxrule=0.45pt,
  left=1.1mm,
  right=1.1mm,
  top=0.9mm,
  bottom=0.9mm
}

\newtcolorbox{greenbox}[1]{
  colback=green!6,
  colframe=success,
  title=\textbf{#1},
  arc=1.5mm,
  boxrule=0.45pt,
  left=1.1mm,
  right=1.1mm,
  top=0.9mm,
  bottom=0.9mm
}

\newcommand{\placeholder}[1]{\textcolor{accent}{\texttt{[#1]}}}
\newcommand{\janpro}{JAN-PRO}

\begin{document}

\begin{center}
  {\LARGE\bfseries Meeting with a \janpro{} Director of Franchise Sales}\\[0.35em]
  {\large Candidate Positioning, Qualification, and Due-Diligence Handout}\\[0.4em]
  {\small Prepared for a commercial-cleaning franchise discovery conversation}
\end{center}

\vspace{0.5em}

\begin{keybox}{Core meeting objective}
The meeting is a two-way qualification conversation. Your job is to present yourself as a financially qualified, operationally serious candidate for a \janpro{} commercial-cleaning franchise while testing whether the brand's account-acquisition model, local support structure, training, territory rules, recurring-revenue assumptions, and service-quality expectations fit your goals and risk profile.
\end{keybox}

\begin{warningbox}{Brand-specific validation note}
Do not rely on informal earnings, account-volume, territory, or startup-cost statements unless they are backed by the current Franchise Disclosure Document (FDD), written franchisor materials, or direct franchisee validation. The FTC requires franchisors to provide an FDD with 23 specific disclosure items, and the FTC states that a prospective franchisee must receive the FDD at least 14 days before being asked to sign a contract or pay money to the franchisor or an affiliate.\cite{ftc-rule,ftc-consumer-guide,ftc-deep-dive}
\end{warningbox}

\section{Positioning Yourself for \janpro{}}

\janpro{} should be approached as a business-to-business services opportunity, not simply as a cleaning job. The strongest candidate message is that you understand the operating model: selling and retaining commercial accounts, hiring or managing reliable labor, meeting service specifications, controlling supplies and equipment costs, resolving customer complaints quickly, and following a standardized franchise system.

\begin{scriptbox}{Opening script tailored to \janpro{}}
``Thank you for taking the time to speak with me. I am evaluating franchise opportunities where my background, capital position, and long-term ownership goals align with a repeatable operating system. I am interested in \janpro{} because commercial cleaning is a recurring B2B service model where account retention, disciplined operations, quality control, and local relationship management matter. My goal today is to understand whether I am a strong candidate for your system and whether your system is the right fit for me.''
\end{scriptbox}

\begin{scriptbox}{Candidate positioning}
``My background is in \placeholder{professional background}. I am comfortable with structured processes, customer service, compliance, technology, operations, and performance measurement. I am looking for an ownership model where I can be an engaged \placeholder{owner-operator / executive operator / multi-unit operator} and build a disciplined commercial-services business in the \placeholder{target market} market.''
\end{scriptbox}

\subsection{Your four-sentence introduction}

\begin{enumerate}
  \item \textbf{Background:} ``My professional background is \placeholder{your background}, with strengths in \placeholder{operations, sales, customer service, hiring, process management, compliance, finance, technology}.''
  \item \textbf{Capital posture:} ``I am evaluating opportunities that fit a liquid-capital range of \placeholder{amount} and a total investment range of \placeholder{amount}.''
  \item \textbf{Ownership goal:} ``My goal is to operate as \placeholder{owner-operator / executive operator / multi-unit operator} and build toward \placeholder{single territory / multiple accounts / multi-unit growth}.''
  \item \textbf{Diligence standard:} ``I plan to review the FDD, speak with current and former franchisees, validate the local market, and build a conservative pro forma before making any commitment.''
\end{enumerate}

\section{What the Director Needs to Learn Quickly}

A franchise sales director will likely evaluate whether you are serious, qualified, coachable, and operationally realistic. For \janpro{}, emphasize that you understand commercial cleaning as a recurring-service business with sales, staffing, training, and quality-control demands.

\begin{longtable}{>{\raggedright\arraybackslash}p{0.28\linewidth}>{\raggedright\arraybackslash}p{0.66\linewidth}}
\toprule
\textbf{Topic} & \textbf{What to communicate} \\
\midrule
Ownership model & Whether you want to be an active owner-operator, executive operator, or growth-oriented owner who may eventually manage teams and multiple accounts. \\
Liquid capital & A clear range, not a vague statement. Example: ``I am evaluating brands that fit a liquid-capital range of \placeholder{amount}.'' \\
Net worth and financing & Whether you expect to self-fund, use SBA-backed financing, bring partners, or combine financing sources. Confirm lender familiarity with commercial-cleaning franchises. \\
Market focus & Target geography, territory interest, and why the local B2B market is attractive: office density, medical/professional space, schools, religious facilities, retail, light industrial, or other commercial demand. \\
Sales readiness & Comfort with prospecting, local networking, commercial proposals, customer follow-up, account onboarding, and renewal/retention discussions. \\
Operations readiness & Ability to manage labor reliability, schedule coverage, supplies, equipment, inspections, complaint resolution, and service consistency. \\
Due-diligence posture & Serious but disciplined: FDD review, franchisee validation, local market analysis, total startup-cost review, and professional legal/accounting review. \\
Timeline & Whether you are exploring, actively qualifying, or ready to move through application, FDD review, validation, discovery, and final decision. \\
\bottomrule
\end{longtable}

\section{Regulatory and Diligence Baseline}

\begin{itemize}
  \item The Federal Trade Commission's Franchise Rule requires franchisors to provide prospective franchisees with a disclosure document containing 23 specific items of information about the franchise offer, officers, and other franchisees.\cite{ftc-rule}
  \item The FTC states that prospective franchise buyers must receive the FDD at least 14 days before they are asked to sign a contract or pay money to the franchisor or its affiliate.\cite{ftc-consumer-guide,ftc-deep-dive}
  \item The FTC advises prospective franchise buyers to read the 23 FDD items, ask questions, evaluate costs, review training and support, speak with franchisees, and treat financial-performance claims carefully.\cite{ftc-consumer-guide,ftc-deep-dive}
  \item SBA-backed financing requires brand eligibility review. SBA states that the SBA Franchise Directory contains franchises and other brands eligible for SBA financial assistance and that placement is not an endorsement or guarantee of success.\cite{sba-directory}
  \item Independent validation matters. The International Franchise Association recommends analyzing the FDD, speaking with current and former franchisees, evaluating initial and ongoing costs, researching market demand, and using independent professional advisors.\cite{ifa-diligence}
\end{itemize}

\section{Financial Qualification Statement}

Be direct about your readiness while avoiding unnecessary personal detail. For a commercial-cleaning franchise, the financial conversation should include not only the franchise fee and startup cost but also working capital, insurance, payroll exposure, vehicle/equipment needs, supplies, local marketing, taxes, and the time required to stabilize recurring accounts.

\begin{scriptbox}{Capital and readiness statement}
``I understand that franchise qualification depends on liquid capital, net worth, financing readiness, and the full initial investment range. I am prepared to discuss my liquid-capital range of \placeholder{amount}, my funding strategy of \placeholder{self-fund / SBA loan / partner capital / combination}, and my expected timeline of \placeholder{timeframe}. I also want to understand the working-capital cushion needed for a commercial-cleaning operation before account revenue becomes predictable.''
\end{scriptbox}

\subsection{Numbers to know before the call}

\begin{longtable}{>{\raggedright\arraybackslash}p{0.39\linewidth}>{\raggedright\arraybackslash}p{0.55\linewidth}}
\toprule
\textbf{Metric} & \textbf{Your prepared answer} \\
\midrule
Liquid capital available & \placeholder{cash, marketable securities, immediately available funds} \\
Estimated net worth & \placeholder{range} \\
Financing plan & \placeholder{cash / SBA / HELOC / investor / retirement rollover / other} \\
Target investment range & \placeholder{maximum all-in startup investment} \\
Working-capital cushion & \placeholder{months of operating cushion beyond initial investment} \\
Target market & \placeholder{city, county, region, or territory} \\
Ownership model & \placeholder{owner-operator, executive operator, multi-account operator} \\
Local customer focus & \placeholder{offices, medical/professional, retail, churches, schools, light industrial, other commercial facilities} \\
Decision timeline & \placeholder{exploring, 3--6 months, 6--12 months, ready now} \\
\bottomrule
\end{longtable}

\section{How to Explain Why \janpro{} Fits}

Avoid saying only that you ``like the brand.'' Tie your interest to operating logic.

\begin{scriptbox}{Brand-fit statement}
``What attracts me to \janpro{} is the potential to build a recurring B2B service business around disciplined execution. I am not looking for a speculative concept. I am looking for a system with clear training, defined standards, support for commercial account growth, measurable quality control, and a franchisee profile that matches how I intend to operate.''
\end{scriptbox}

\begin{scriptbox}{More advanced version}
``I understand that commercial cleaning depends on sales discipline, customer retention, labor reliability, inspection standards, and cost control. I want to understand how \janpro{} helps franchisees acquire accounts, price work, train staff, resolve service issues, retain customers, and scale without letting quality slip.''
\end{scriptbox}

\section{Strategic Questions to Ask \janpro{}}

Ask questions that sound like an investor and operator, not a casual buyer. The goal is to test the franchise structure, account-acquisition model, economics, support, territory, and candidate fit.

\subsection{Franchise structure and local support}

\begin{itemize}
  \item ``Who would be my direct franchisor or support organization in this market: the national franchisor, a regional developer, or another local franchise entity?''
  \item ``Which party provides training, account support, billing support, field inspections, complaint escalation, and ongoing coaching?''
  \item ``How is the relationship between national brand standards and local/regional support managed?''
  \item ``What should I understand about the difference between a unit franchise, regional franchise, master franchise, or development opportunity if those structures apply in this market?''
  \item ``What are the most common misunderstandings candidates have about how the \janpro{} system works?''
\end{itemize}

\subsection{Account acquisition, leads, and sales expectations}

\begin{itemize}
  \item ``How are commercial accounts generated: franchisor-provided leads, regional sales support, owner prospecting, referrals, local marketing, or a combination?''
  \item ``Does the system provide initial accounts, guaranteed account volume, account packages, or lead targets? If so, where is that documented in the FDD or agreement?''
  \item ``What level of local sales activity do successful franchisees perform themselves?''
  \item ``What are typical sales-cycle realities for commercial cleaning accounts in this market?''
  \item ``How does \janpro{} help franchisees prepare proposals, price jobs, define scope, and avoid underbidding?''
  \item ``What happens if a customer cancels, reduces scope, delays payment, or disputes service quality?''
\end{itemize}

\subsection{Territory and customer concentration}

\begin{itemize}
  \item ``What exactly is protected: geography, customer accounts, account type, lead source, or something else?''
  \item ``Can another \janpro{} franchisee serve accounts in or near my market? Under what circumstances?''
  \item ``How are national, regional, and local commercial accounts assigned?''
  \item ``What customer-concentration risk should I plan for if early revenue depends on a small number of accounts?''
  \item ``Which FDD item and agreement sections define territory, account ownership, non-compete, renewal, transfer, and termination rights?''
\end{itemize}

\subsection{Training, quality control, and operations}

\begin{itemize}
  \item ``What does onboarding look like from signing through first account launch?''
  \item ``How much training is technical cleaning training versus sales, customer management, labor management, finance, and quality control?''
  \item ``How are cleaning standards, inspections, customer complaints, and corrective actions documented?''
  \item ``What service failures most commonly hurt new franchisees, and how does \janpro{} coach owners through them?''
  \item ``What software, checklists, reporting tools, or inspection methods are required?''
  \item ``What are the expectations for night, weekend, holiday, backup, and emergency coverage?''
\end{itemize}

\subsection{Labor, compliance, insurance, and safety}

\begin{itemize}
  \item ``What labor model do your successful franchisees use: owner-performed work, employees, subcontractors where permitted, or a mix?''
  \item ``What insurance, bonding, background-check, workers' compensation, licensing, or local business requirements are typical in this market?''
  \item ``How are chemical handling, safety data sheets, equipment training, OSHA expectations, and customer-site safety rules handled?''
  \item ``Are there special requirements for medical, educational, government, industrial, or regulated facilities?''
  \item ``What costs do new franchisees commonly underestimate for labor, payroll taxes, insurance, supplies, equipment replacement, vehicle use, and supervision?''
\end{itemize}

\subsection{FDD, economics, and financing}

\begin{itemize}
  \item ``Does the current FDD include Item 19 financial performance representations, and how should I interpret them conservatively?''
  \item ``Which revenue, gross margin, account-retention, cancellation, and ramp-up assumptions should I validate through franchisee calls?''
  \item ``What are the ongoing fees: royalties, brand/marketing fund, account administration, technology, training, renewal, transfer, insurance, or other required payments?''
  \item ``What is the minimum working-capital cushion you recommend beyond the stated initial investment range?''
  \item ``Do franchisees commonly use SBA-backed financing, and is the brand currently eligible in the SBA Franchise Directory?''
  \item ``Can you introduce me to lenders who understand the commercial-cleaning franchise model, without pressuring me toward one financing source?''
\end{itemize}

\section{Jan-Pro Candidate Fit Narrative}

Use this narrative if the Director asks, ``Why \janpro{}?''

\begin{keybox}{Polished answer}
``I am drawn to \janpro{} because I want a service business with recurring commercial demand, operational discipline, and measurable quality standards. My strengths are process management, customer follow-up, technology, and structured execution. I understand that success will depend on local sales, account retention, labor reliability, cost control, and following the system. I am not looking for passive income. I am looking for a franchise system where I can operate professionally, validate the economics carefully, and scale only after the first accounts are stable.''
\end{keybox}

\section{What Not to Say}

\begin{longtable}{>{\raggedright\arraybackslash}p{0.38\linewidth}>{\raggedright\arraybackslash}p{0.56\linewidth}}
\toprule
\textbf{Avoid saying} & \textbf{Say this instead} \\
\midrule
``How much money can I make?'' & ``What financial performance representations are disclosed in Item 19, and what assumptions should I use when building a conservative pro forma?'' \\
``Is this passive income?'' & ``What level of owner involvement do your strongest franchisees typically have during the first year?'' \\
``Will you give me accounts?'' & ``How are accounts generated, assigned, priced, onboarded, retained, and documented in the FDD or franchise agreement?'' \\
``Cleaning is simple, so this should be easy.'' & ``I understand that commercial cleaning requires sales discipline, labor reliability, quality control, customer retention, and cost management.'' \\
``Can I negotiate the franchise agreement?'' & ``Which terms are standard across the system, and which items may vary by market, account structure, or development path?'' \\
``I do not know my budget yet.'' & ``I am targeting an all-in investment range of \placeholder{amount} and want to confirm whether that aligns with your current requirements.'' \\
``I just want the cheapest route to ownership.'' & ``I want to understand the investment level that gives a new franchisee a realistic operating runway and adequate working capital.'' \\
\bottomrule
\end{longtable}

\section{FDD Review Focus for a \janpro{} Candidate}

\begin{longtable}{>{\raggedright\arraybackslash}p{0.22\linewidth}>{\raggedright\arraybackslash}p{0.72\linewidth}}
\toprule
\textbf{FDD item} & \textbf{Why it matters for \janpro{} diligence} \\
\midrule
Items 1--2 & Understand the franchisor, affiliates, regional structure, executives, and operating history. Confirm who actually supports you locally. \\
Items 3--4 & Review litigation and bankruptcy history for franchisor, affiliates, and key executives. Look for patterns involving franchisee relationships, fees, accounts, termination, or support. \\
Items 5--7 & Build the full startup and first-year cash model: franchise fee, equipment, supplies, insurance, software, training, vehicle, payroll, local marketing, professional fees, and working capital. \\
Item 8 & Understand required suppliers, cleaning chemicals, equipment, uniforms, software, insurance, and whether the franchisor receives rebates or financial benefits from designated vendors. \\
Item 11 & Validate training, operations manuals, account support, marketing, sales assistance, technology, inspections, and field support. Ask recent franchisees whether the support matched expectations. \\
Item 12 & Read territory, account, lead, customer, and online-marketing restrictions carefully. Territory language is critical in a commercial-services model. \\
Item 17 & Review renewal, termination, transfer, dispute resolution, non-compete, and post-termination obligations before signing. \\
Item 19 & Treat financial performance representations conservatively. Ask what population of franchisees the numbers represent and whether the data reflects your target market and operating model. \\
Item 20 & Use the contact lists for current and former franchisee validation calls. Ask about account acquisition, ramp-up, support quality, cancellations, costs, and whether they would buy again. \\
Item 21 & Have an accountant review financial statements to assess whether the franchisor can deliver promised support. \\
\bottomrule
\end{longtable}

\section{Franchisee Validation Call Sheet}

Speak with current and former franchisees before making a final decision. The FTC and IFA both emphasize franchisee conversations as a key source of practical diligence.\cite{ftc-consumer-guide,ftc-deep-dive,ifa-diligence}

\begin{enumerate}
  \item How long did it take from signing to launching your first accounts?
  \item What did your actual startup cost and working-capital need look like versus the FDD range?
  \item How were your initial accounts generated, and what did you personally have to do to grow revenue?
  \item What surprised you most about labor, scheduling, inspections, supplies, insurance, or customer complaints?
  \item How responsive was local/regional support during the first 90, 180, and 365 days?
  \item What accounts were hardest to service profitably?
  \item What fees or costs did you underestimate?
  \item What is your customer-retention experience, and why do customers cancel?
  \item How many hours per week did you work in year one compared with what you expected?
  \item Knowing what you know now, would you buy this franchise again? Why or why not?
\end{enumerate}

\section{Closing the Meeting}

End by forcing clarity on fit and next steps.

\begin{scriptbox}{Professional close}
``This has been helpful. Based on what we discussed, I would like to understand the next stage in your process. What materials should I review, what qualification information do you need from me, and what is the expected timeline from this initial conversation to application, FDD review, validation calls, discovery, and a possible franchise agreement?''
\end{scriptbox}

\begin{scriptbox}{Direct fit question}
``Based on what I have shared, do you see me as a strong candidate for \janpro{}? Are there any concerns about my financial profile, market, owner involvement, sales readiness, or operating fit that I should address before moving forward?''
\end{scriptbox}

\section{Post-Meeting Follow-Up Email Template}

\begin{scriptbox}{Follow-up email}
\textbf{Subject:} Follow-up on \janpro{} franchise discussion\\[0.4em]
\textbf{Body:}\\
Hello \placeholder{Name},\\[0.25em]
Thank you for taking the time to speak with me about the \janpro{} franchise opportunity. I appreciated the overview of the commercial-cleaning model, candidate profile, training process, account-development expectations, local support structure, and market opportunity.\\[0.25em]
Based on our conversation, I remain interested in evaluating the opportunity further. My next diligence priorities are reviewing the current FDD, understanding the full investment range and working-capital expectations, validating the account-acquisition and support model, reviewing territory/account rules, and speaking with current and former franchisees where appropriate.\\[0.25em]
Please send the next-step materials and any qualification documents you need from me. I am also happy to provide additional detail on my liquid capital, financing plan, target market, ownership timeline, and operating goals.\\[0.25em]
Best,\\
\placeholder{Your Name}
\end{scriptbox}

\section{Pre-Call Checklist}

\begin{itemize}
  \item Confirm your liquid-capital range and maximum all-in investment range.
  \item Decide whether your target role is owner-operator, executive operator, or future multi-account operator.
  \item Identify your target market and why it makes sense for B2B commercial cleaning.
  \item Prepare a local market thesis: commercial density, office/medical/professional corridors, local competition, commute radius, labor availability, and target account types.
  \item Review \janpro{} public materials and prepare questions that cannot be answered by the website.
  \item Prepare a question set around account generation, assignment, retention, cancellation, and pricing.
  \item Confirm whether SBA-backed financing is part of your plan and check the SBA Franchise Directory before relying on SBA eligibility.
  \item Prepare a post-meeting diligence plan: FDD review, franchisee calls, professional review, financing review, market review, and conservative pro forma.
\end{itemize}

\section{Red Flags to Watch For}

\begin{warningbox}{Warning signs during the sales conversation}
\begin{itemize}
  \item Pressure to pay quickly before you have reviewed the FDD and agreements.
  \item Informal earnings or account-volume claims that are not tied to Item 19 or written materials.
  \item Vague answers about how accounts are generated, assigned, priced, serviced, retained, or replaced after cancellation.
  \item Overemphasis on ``recurring revenue'' without clear discussion of labor, cancellations, service failures, and customer concentration.
  \item Resistance to current and former franchisee validation calls.
  \item Unclear distinction between national franchisor support and local/regional support.
  \item No direct answer on territory, account ownership, competition among franchisees, or national-account assignment.
  \item Dismissive treatment of insurance, labor law, safety, chemical handling, or site-specific compliance requirements.
\end{itemize}
\end{warningbox}

\section{Best Overall Script}

\begin{greenbox}{Use this when you need a polished summary}
``I am looking for a franchise opportunity where my capital, professional background, and ownership goals align with a proven operating system. I am interested in \janpro{} because commercial cleaning appears to reward disciplined B2B sales, account retention, quality control, and operational consistency. I want to evaluate the opportunity carefully through the FDD, franchisee validation, unit economics, training model, local support structure, and territory/account rules. My goal today is to understand whether I am a strong candidate for your system and whether the system is the right fit for me.''
\end{greenbox}

\renewcommand{\refname}{\color{accent}References}
\begin{thebibliography}{9}

\bibitem{ftc-rule}
Federal Trade Commission.
\textit{Franchise Rule}.
\url{https://www.ftc.gov/legal-library/browse/rules/franchise-rule}

\bibitem{ftc-consumer-guide}
Federal Trade Commission.
\textit{A Consumer's Guide to Buying a Franchise}.
\url{https://www.ftc.gov/business-guidance/resources/consumers-guide-buying-franchise}

\bibitem{ftc-deep-dive}
Federal Trade Commission.
\textit{Franchise Fundamentals: Taking a Deep Dive into the Franchise Disclosure Document}.
\url{https://www.ftc.gov/business-guidance/blog/2023/05/franchise-fundamentals-taking-deep-dive-franchise-disclosure-document}

\bibitem{sba-directory}
U.S. Small Business Administration.
\textit{SBA Franchise Directory}.
\url{https://www.sba.gov/business-guide/plan-your-business/buy-existing-business-or-franchise/sba-franchise-directory}

\bibitem{ifa-diligence}
International Franchise Association.
\textit{Diligence in Franchising}.
\url{https://www.franchise.org/franchising-overview/diligence-in-franchising/}

\end{thebibliography}


\end{document}
